\section{Paso 6: Conclusiones Individuales y Declaración de Uso de IA}
\subsection{Conclusiones (Estructura Tripartita: Afirmación, Evidencia Normativa e Implicación en el PFC)}
\begin{enumerate}
	\item \textbf{Conclusión 1 --- Gilces Carranza José Ignacio (Sobre la Especificación EARS y la Propiedad ``Apropiado al Nivel'' en IEEE 29148):} \\
	\textbf{(1) Afirmación:} La redacción de requisitos mediante patrones sintácticos estructurados y la abstención estricta de incluir tecnologías de implementación en la especificación eliminan la ambigüedad semántica y preservan la independencia entre el problema de negocio y el diseño posterior. \textbf{(2) Evidencia con cita:} Conforme lo establecen la norma ISO/IEC/IEEE 29148:2018 (cláusula 5.2.5) \cite{ieee29148} y la metodología EARS de Mavin et al. \cite{mavin2009ears}, todo requisito individual bien formado debe ser inambiguo, singular, verificable y apropiado al nivel (describiendo el ``qué'' necesita el stakeholder sin imponer prematuramente el ``cómo'' tecnológico). \textbf{(3) Implicación para el PFC y la profesión:} En nuestra Plataforma de Ventas para almacenes eléctricos de Quevedo, depurar los 15 requisitos de menciones a herramientas de programación específicas y expresarlos mediante umbrales observables ($\le 3.0\text{ s}$ en recomendaciones técnicas filtradas por existencias reales $>0$) nos permitió construir una especificación limpia, verificable y profesionalmente auditable.
	
	\item \textbf{Conclusión 2 --- Gilces Carranza José Ignacio (Sobre la Cuantificación de Atributos de Calidad con ISO/IEC 25010:2023):} \\
	\textbf{(1) Afirmación:} La viabilidad operativa de un sistema transaccional acoplado a asistencia inteligente depende decisivamente de la cuantificación rigurosa de sus Requisitos No Funcionales (RNF) en términos de eficiencia de desempeño, fiabilidad, seguridad y capacidad de interacción. \textbf{(2) Evidencia con cita:} De acuerdo con el estándar ISO/IEC 25010:2023 \cite{iso25010} y la clasificación de Glinz \cite{glinz2007nfr}, los atributos de calidad acotan el comportamiento del producto bajo condiciones reales de operación y definen criterios objetivos de aceptación. \textbf{(3) Implicación para el PFC y la profesión:} Vincular cada RNF del PFC al modelo ISO/IEC 25010 nos permitió exigir una disponibilidad del $\ge 99.5\%$, un mecanismo de contingencia local en $\le 200\text{ ms}$ ante interrupciones del servicio externo de IA y estándares de accesibilidad visual WCAG 2.1 AA con operación íntegra por teclado, garantizando continuidad absoluta en el cobro en horas pico.
	
	\item \textbf{Conclusión 3 --- Vélez Vinces Jostin Gustavo (Sobre el Límite del Sistema y las Ocho Perspectivas de Contexto de Pohl):} \\
	\textbf{(1) Afirmación:} Explorar sistemáticamente el entorno sociotécnico donde operará el sistema mediante perspectivas multidimensionales permite descubrir restricciones físicas, culturales y regulatorias críticas que permanecerían ocultas en un levantamiento convencional. \textbf{(2) Evidencia con cita:} Según Klaus Pohl (2025) \cite{pohl2025} y el sílabo IREB CPRE v3.1 \cite{ireb2023}, el contexto del sistema debe analizarse a través de ocho perspectivas complementarias (sujeto, uso, entorno TI, desarrollo, física, social, legal y económica) para delimitar con exactitud la zona gris y las interfaces externas. \textbf{(3) Implicación para el PFC y la profesión:} Aplicar las ocho perspectivas al comercio eléctrico de Quevedo hizo emerger requisitos esenciales para el PFC, tales como el reconocimiento de jerga ferretera local (``breaker'', ``taipe''), la prohibición de enviar datos personales de compradores hacia servicios externos en cumplimiento de la LOPDP y la integración con el esquema de comprobantes electrónicos del SRI.
	
	\item \textbf{Conclusión 4 --- Robinson Espinoza Jeanpierre (Sobre el Proceso Iterativo de IR y la Negociación de Conflictos entre Stakeholders):} \\
	\textbf{(1) Afirmación:} La ingeniería de requisitos constituye un proceso iterativo y continuo de elicitación, análisis, negociación, especificación y validación cuyo propósito central es conciliar intereses contrapuestos de múltiples actores en una especificación consistente. \textbf{(2) Evidencia con cita:} Tal como sostienen el SWEBOK v4.0 \cite{swebok2024}, el PMBOK 7.ª edición \cite{pmbok} y Wiegers \& Beatty \cite{wiegers2013}, priorizar a los interesados mediante una matriz de poder/interés y resolver conflictos de alcance en etapas tempranas previene el retrabajo costoso en fases posteriores del ciclo de vida. \textbf{(3) Implicación para el PFC y la profesión:} En el PFC, la fase de análisis nos permitió resolver la tensión entre el propietario (quien restringía el gasto en equipamiento local) y el cajero (quien demandaba sugerencias técnicas en $\le 3.0\text{ s}$) mediante la especificación de una interfaz hacia un servicio externo de inferencia IA combinada con reglas de filtrado sobre el inventario disponible.
	
	\item \textbf{Conclusión 5 --- Robinson Espinoza Jeanpierre (Sobre la Trazabilidad Bidireccional y el Rigor Documental en la SRS):} \\
	\textbf{(1) Afirmación:} Asegurar la trazabilidad bidireccional desde las fuentes primarias y los stakeholders de origen hasta los artefactos de especificación y validación es un requisito indispensable para mantener la integridad del proyecto a lo largo de su ciclo de vida. \textbf{(2) Evidencia con cita:} Las normas ISO/IEC/IEEE 29148:2018 \cite{ieee29148}, ISO/IEC/IEEE 12207:2017 \cite{ieee12207} e ISO/IEC/IEEE 15288:2023 \cite{ieee15288} determinan que un conjunto de requisitos solo es de calidad si es internamente consistente, completo y trazable tanto hacia atrás (necesidad de negocio) como hacia adelante (criterios de verificación). \textbf{(3) Implicación para el PFC y la profesión:} La construcción articulada de los artefactos de elicitación (\texttt{ACTA-ELI-01}), análisis (\texttt{ANA-01}), especificación (\texttt{ESP-01}) y validación (\texttt{VAL-01}) nos proporciona una línea base formal y verificable para las siguientes etapas de la carrera de Ingeniería de Software.
\end{enumerate}
\subsection{Declaración Explícita de Uso Ético y Validación Crítica de Inteligencia Artificial}
En estricto cumplimiento de la Sección 4 y del Paso 6d de la Guía de Práctica Experimental de la UTEQ, los autores (\textbf{Gilces Carranza José Ignacio, Vélez Vinces Jostin Gustavo y Robinson Espinoza Jeanpierre}) declaran de manera transparente que durante la elaboración del presente informe se utilizaron herramientas de Inteligencia Artificial Generativa como apoyo de asistencia en la estructuración tipográfica, revisión de estilo técnico y verificación sintáctica de plantillas EARS. Todo el contenido conceptual, la selección y contextualización del sistema real en Quevedo (PFC), las métricas cuantitativas de los 15 requisitos (verificando que se mantengan libres de decisiones prematuras de implementación), el análisis de las ocho perspectivas de contexto y las citas bibliográficas fueron contrastados, verificados y validados críticamente por los integrantes del equipo contra las fuentes primarias obligatorias (ISO/IEC/IEEE 29148:2018 \cite{ieee29148}, ISO/IEC 25010:2023 \cite{iso25010}, Pohl 2025 \cite{pohl2025} y SWEBOK v4.0 \cite{swebok2024}), asumiendo los autores la responsabilidad académica e intelectual íntegra sobre el documento final.
\section{Referencias Bibliográficas}
\printbibliography[heading=none]
\newpage
\section*{Anexos Documentales (Paso 2c y 2d)}
\subsection*{Anexo A: Exportación Estandarizada de Fuentes Primarias en Formato IEEE Comprimido}
A continuación se detalla el listado maestro exportado desde la colección grupal con las 12 fuentes primarias normativas y bibliográficas empleadas y citadas a lo largo del documento:
\begin{enumerate}
	\item \textbf{[1] Norma Principal:} ISO/IEC/IEEE, \textit{ISO/IEC/IEEE 29148:2018 - Systems and software engineering -- Life cycle processes -- Requirements engineering}, IEEE, 2018. DOI: \texttt{10.1109/IEEESTD.2018.8559686}.
	\item \textbf{[2] Ciclo de Vida de Sistemas:} ISO/IEC/IEEE, \textit{ISO/IEC/IEEE 15288:2023 - Systems and software engineering -- System life cycle processes}, IEEE, 2023.
	\item \textbf{[3] Ciclo de Vida de Software:} ISO/IEC/IEEE, \textit{ISO/IEC/IEEE 12207:2017 - Systems and software engineering -- Software life cycle processes}, IEEE, 2017.
	\item \textbf{[4] Modelo de Calidad SQuaRE:} International Organization for Standardization, \textit{ISO/IEC 25010:2023 - Systems and software quality requirements and evaluation (SQuaRE) -- Product quality model}, ISO, 2023.
	\item \textbf{[5] Texto Base de IR:} K. Pohl, \textit{Requirements Engineering: Fundamentals, Principles, and Techniques}, 2.ª ed. Springer, 2025. ISBN: \texttt{978-3-662-69204-2}.
	\item \textbf{[6] Cuerpo de Conocimiento:} H. Washizaki (Ed.), \textit{Guide to the Software Engineering Body of Knowledge (SWEBOK Guide) v4.0}, IEEE Computer Society, 2024.
	\item \textbf{[7] Gestión de Proyectos y Stakeholders:} Project Management Institute, \textit{A Guide to the Project Management Body of Knowledge (PMBOK Guide)}, 7.ª ed. PMI, 2021.
	\item \textbf{[8] Certificación Internacional IREB:} International Requirements Engineering Board, \textit{CPRE Foundation Level Syllabus v3.1}, IREB, 2023.
	\item \textbf{[9] Declaración de Visión:} G. A. Moore, \textit{Crossing the Chasm: Marketing and Selling Disruptive Products to Mainstream Customers}, 3.ª ed. HarperBusiness, 2014.
	\item \textbf{[10] Sintaxis Estructurada EARS:} A. Mavin, P. Wilkinson, A. Harwood y M. Novak, ``Easy Approach to Requirements Syntax (EARS),'' en \textit{Proc. 17th IEEE Int. Requirements Engineering Conf. (RE'09)}, 2009, pp. 317--322.
	\item \textbf{[11] Requisitos No Funcionales:} M. Glinz, ``On Non-Functional Requirements,'' en \textit{Proc. 15th IEEE Int. Requirements Engineering Conf. (RE'07)}, 2007, pp. 21--26.
	\item \textbf{[12] Ingeniería de Requisitos Aplicada:} K. Wiegers y J. Beatty, \textit{Software Requirements}, 3.ª ed. Microsoft Press, 2013; y E. Hull, K. Jackson y J. Dick, \textit{Requirements Engineering}, 3.ª ed. Springer, 2011.
\end{enumerate}
\subsection*{Anexo B: Evidencia de Gestión Bibliográfica (Captura de Biblioteca Grupal)}
En la Figura \ref{fig:anexo_zotero} se adjunta la captura de pantalla de la biblioteca grupal configurada en el gestor bibliográfico, evidenciando la organización y clasificación de las fuentes primarias obligatorias de la Unidad I (ver Anexo B).
\begin{figure}[htbp]
	\centering
	\includegraphics[width=0.95\textwidth]{figura/zotero_evidencia.png}
	\caption{Anexo B: Captura de pantalla de la colección bibliográfica grupal con las fuentes primarias en formato IEEE.}
	\label{fig:anexo_zotero}
\end{figure}
